% !TeX spellcheck = en_US
\documentclass[french]{yLectureNote}

\title{Électromagnétisme 2 PS}
\subtitle{Physique}
\author{Paulhenry Saux}
\date{\today}
\yLanguage{Français}
%lionels.calmels@cemes.fr
\professor{M.Brut}
\usepackage{graphicx}%----pour mettre des images
\usepackage[utf8]{inputenc}%---encodage
\usepackage{geometry}%---pour modifier les tailles et mettre a4paper
%\usepackage{awesomebox}%---pour les boites d'exercices, de pbq et de croquis ---d\'esactiv\'e pour les TP de PC
\usepackage{tikz}%---pour deiffner + d\'ependance de chemfig
% \usepackage{tabularx}%---pour dimensionner automatiquement les tableaux avec variable X
\usepackage{awesomebox}%---Pour les boites info, danger et autres
\usepackage{menukeys}%---Pour deiffner les touches de Calculatrice
\usepackage{fancyhdr}%---pour les en-t\^ete personnalis\'ees
\usepackage{blindtext}%---pour les liens
\usepackage{hyperref}%---pour les liens (\`a mettre en dernier)
\usepackage{caption}%---pour la francisation de la l\'egende table vers Tableau
\usepackage{pifont}
\usepackage{array}%---pour les tableaux
\usepackage{lipsum}
\usepackage{yFlatTable}
\usepackage{multicol}
\newcommand{\Lim}[1]{\lim\limits_{\substack{#1}}\:}
\renewcommand{\vec}{\overrightarrow}
\newcommand{\N}[0]{\mathbb{N}}
\newcommand{\dd}{\mathrm{d}}
\newcommand{\norm}[1]{||\vec{#1}||}
\newcommand{\fo}{\psi(\vec{r},t)}
\newcommand{\foe}{\psi(\vec{r},t)\*}
\newcommand{\HH}{\hat{H}}
\newcommand{\hb}{\hbar}
\newcommand{\lap}{\nabla^2}
\newcommand{\lapcc}{\frac{\partial^2 }{\partial x^2}+\frac{\partial^2 }{\partial y^2}+\frac{\partial^2 }{\partial z^2}}
\newcommand{\E}{\vec{E}(M)}
\newcommand{\B}{\vec{B}(M)}
\newcommand{\V}{V(M)}
\newcommand{\er}{\vec{e_{\rho}}}
\newcommand{\ep}{\vec{e_{\varphi}}}
\newcommand{\pip}{\pi^+}
\newcommand{\pim}{\pi^-}
\newcommand{\mv}{\mathcal{V}}
\DeclareMathOperator{\Div}{Div}
\newcommand{\rot}{\vec{\mathrm{rot}}}
\newcommand{\grad}{\vec{\mathrm{grad}}}
\setcounter{chapter}{1}
\begin{document}
\chapter{Induction}
\section{Approche expérimentale}
En rentrant un aimant dans une bobine\marginCritical{Ce n'est pas que la présence de l'aimant, il faut qu'il y ait un déplacement de ce dernier !}, on observe l'apparition d'un courant. C'est un courant induit.

Le signe du courant s'inverse si je bouge l'aimant dans l'autre sens. Son intensité dépend de 
\begin{itemize}
    \item la vitesse de l'aimant
    \item l'amplitude du déplacement
    \item l'intensité de B
    \item la géométrie de la bobine
    \item etc
\end{itemize}
Ce n'est pas B qui crée le courant, mais sa variation.

Si j'utilise un voltmètre aux bornes de la bobine, je mesure une tension. C'est la tension générateur appelée force électromotrice (fem).

\subsection{Faire varier B}
\begin{itemize}
    \item En déplaçant l'aimant ou la bobine
    \item En modifiant l'intensité de B
    \item En modifiant la surface de la spire ou sont orientation
\end{itemize}
\subsection{Force électromotrice}
\begin{definition}[Force électromotrice (fem)]
    Toute variation de flux magnétique à travers un circuit provoque l'apparition d'une fem. Le circuit est alors équivalent à
    une source de tension notée $e$. Si le circuit est ouvert, la tension vaut $e$, dans le cas contraire, elle induit l'apparition d'un courant induit noté $i$.
\end{definition}
\begin{itemize}
    \item Elle est homogène à une tension et s'exprime en V 
    \item $e(t) = \int \E_m\cdot \vec{\dd l}$ où $\E_m$ est le champ électromoteur
\end{itemize}
\section{Lois de l'induction}
\begin{theorem}[Loi de Faraday]
    $$e(t) = -\frac{\dd \phi}{\dd t}$$ avec $\phi = \iint_S \B\cdot \vec{\dd S}$.
\end{theorem}
\begin{theorem}[Loi de Lenz]
    Les phénomènes d'induction s'opposent par leur effets aux causes qui leur ont donné naissance. On en déduit que le sens du
    courant induit est tel que le flux de son propre champ magnétique s'oppose à la variation de flux inducteur.
\end{theorem}
\begin{theorem}[Équation de Maxwell Faraday]
    $$\rot \E = -\frac{\partial \B}{\partial t}$$
\end{theorem}
\begin{myproof}
    Que devient $\rot \E$ en régime variable ?
    \begin{flalign*}
        e(t) &= \int \E\cdot \vec{\dd l}\\
        &= -\frac{\dd \phi}{\dd t}\\
        &= -\frac{\dd }{\dd t}\iint_S \B\cdot \vec{\dd S}\\
        &= -\iint_S \frac{\partial \B}{\partial t}\cdot \vec{\dd S}
    \end{flalign*}
    Avec Stokes, on a $$\int \E\cdot \vec{\dd l} = \iint_S \rot(\E)\cdot \vec{\dd S}$$
Par identification, on obtient $\rot \E = -\frac{\partial \B}{\partial t}$
\end{myproof}
\subsection{Circuit fixe dans un champ variable}
On considère une spire de rayon r, de résistance R dans le plan (Oxy).

On a $\B(t) = B_m \cos(\omega t)\vec{e_z}$. On oriente $\vec{\dd S}$ selon $\vec{e_z}$.

Calculons le flux : 
\begin{flalign*}
    \phi &= \iint_S \B\cdot \vec{\dd S}\\
    &= B(t)\pi r^2\\
    &= B_m \cdot(\omega t) \pi r^2\\
\end{flalign*}

On en déduit $e$ :
\begin{flalign*}
    e&= -\frac{\dd \phi}{\dd t}\\
    &= B_m \omega \pi r^2 \sin(\omega t)
\end{flalign*}
On en déduit $i(t)$ : $$i(t) = \frac{e(t)}{R} = \frac{B_m \omega \pi r^2}{R} \sin(\omega t)$$

\includegraphics[scale=0.5]{Figure_1}

On voit que $i(t)$ est toujours opposé aux taux de variation $\frac{\dd \phi}{\dd t}$, ce qui est 
conforme à la loi de Lenz. %De plus, pour 

S'entrainer avec bobine tourant, en rotation, etc
\subsection{Fem dans un circuit mobile}
On considère un cuircuit C mobile dans un référenriel R où le champ $\E,\B$ est statique.
On considère une charge au repos dans C, elle se déplace à la m\^eme  vitesse ue C ($\vec{v}$).

Elle peut \^etre mise en mouvement par la force de Lorenz, donc 

\explanation{b1}{$\vec{F}$ est la force de Lorenz, en circuit fixe, donc il a quand meme un champ électrique $\vec{F} = q\vec{E}$ qui est le champ électromoteur. 
Dans le cas de l'induction de Lorenz, on aurait $\vec{F} = q(\vec{v}\wedge \vec{B})$. Ici, on a les deux pour la démonstration.}

\begin{flalign*}
    \vec{F} &= q(\E+\vec{v}\wedge \B)\\
    e(t) &= \int \frac{\vec{F}}{q}\vec{\dd l}\explain{b1}{right}{0}{0.5}{}\\
    &= \int_C (\E+\vec{v}\wedge \B)\cdot \vec{\dd l}\\
    &= \int_C \E\cdot \vec{\dd l} + \int_C (\vec{v}\wedge \B)\cdot \vec{\dd l}\\
    &= -\iint_S \frac{\partial \B}{\partial t} \cdot \vec{\dd S}+ \int_C (\vec{v}\wedge \B)\cdot \vec{\dd l}
\end{flalign*}
\begin{theorem}[Fem dans circuit mobile]
    $$-\iint_S \frac{\partial \B}{\partial t} \cdot \vec{\dd S}+ \int_C (\vec{v}\wedge \B)\cdot \vec{\dd l}$$
\end{theorem}
\warningInfo{Types d'induction}{
\begin{itemize}
    \item induction statique (Neumann, $\vec{v}=0$) : on a $e(t) = -\iint \frac{\partial \B}{\partial t}\cdot \vec{\dd S}$
    \item Induction de Lorenz : On a $\frac{\partial \B}{\partial t} = 0$, donc 
     $$e(t) = \int_C (\vec{v}\wedge \B)\cdot \vec{\dd l}$$
\end{itemize}
}
\section{Auto-Induction et inductance d'une bobine}
\subsection{Phénomène d'auto-Induction}
On considère une spire avec une inensité variable. On observe le champ crée par l'inetsnité mais il y a aussi le 
champ induit par la variation du courant. Cela crée un nouveau courant, qui va s'opposer au courant initial. C'est l'auto-induction.

\begin{definition}[Auto-induction]
    La spire parcourue par un courant crée un champ B qui crée un courant induit dans cette spire.
    Par la loi de Lorenz, le courant induit s'oppose au courant initial.
\end{definition}
\subsection{Induction propre et inductance}
\begin{definition}[Flux propre]
    Flux crée par la spire à l'intérieur d'elle-m\^eme, dépend de B, qui dépend des caractéristiques
    de la spire et de l'intensité du courant.
\end{definition}
On peut donc écrire $\phi = L i(t)$, avec L qui est une constante en henry. 

On a donc $$e= -\frac{\dd \phi}{\dd t} = -L\frac{\dd i}{\dd t}$$
\subsection{Inductance mutuelle entre deux circuits}
On prend 2 circuits quelconques appelés $C_1$ et $C_2$ parcourus par des intensités $i_1$ et $i_2$. Le point O est à une certaine distance des 2 circuits, $\vec{r_1}, \vec{r_2}$.

On note $\phi_{21}$ le flux duc champ $\vec{B_1}$ crée à travers $C_2$

Comme 
$$\phi_{21}  = \iint_{S_2} \vec{B_1}\cdot \dd\vec{S_2}$$
on a
$B_1 \propto i_1 \Rightarrow \phi_{21}\propto i_1$. On définit l'inductance $L_{21} = \frac{\phi_{21}}{i_{1}}$ l'inductance mutuelle. Par symétrie, on a 
$$M = L_{21} = L_{12}$$

En effet, 
\begin{flalign*}
    \phi_{21} &= \iint_{S2}\vec{B_1}\cdot\vec{\dd S_2}\\
    &= \int_{C_2}\vec{A_1}\cdot \vec{\dd l_2}\\
    &= \int_{C_2}(\frac{\mu_0i_1}{4\pi}\int_{C_1}\frac{\vec{\dd l_1}}{||\vec{r_2}-\vec{r_1}||})\cdot \vec{\dd l_2}\\
    &= \frac{\mu_0}{4\pi} \int_{C_2}\int_{C1}\frac{\vec{\dd l_1}\cdot \vec{\dd l_2}}{||\vec{r_2}-\vec{r_1}||}i_1
\end{flalign*}
Donc $L_{21}$  ne dépend que de la géométrie des cirucits et de leur orientation relative.
\subsection{Matrices d'inductance}
On se place dans les m\^emes conditions que dans la partie précédente.

On considère toujours deux circuits $C_1$ et $C_2$, on note le flux total à travers $C_1$ : 

\explanation{b2}{Flux dans le cirucit 1 crée par le circuit 2}

$$\phi_1 = \phi_{12}\explain{b2}{right}{0}{0.5}{}+\phi_{11}$$

On obtient pour le circuit 2 : $\phi_{2} = \phi_{21}+\phi_{22}$

Or, $\phi{11} = L_1i_1, \phi_{12} = Mi_2, \phi_{22}=L_2i_2, \phi_{21}M i_1$

On écrit la matrice d'inductance : %TODO Matrice à générer phi = phi1 phi2 en connelone, pareil pour i, et on écrit phi = L i avec L = L1 M//M L2

\checkInfo{Application au cas de 2 bobines en série}{
\begin{flalign*}
    \phi &= \phi_1+\phi_2\\
    &= (L_1+L_2+2M)i
\end{flalign*}
Les inductances ne sont donc pas additives, sauf si $M=0$.
}
\section{Transformateurs}
Principe : 2 bobines sont placées face à face dans 2 circuits différents. l'objectif est de montrer que la différence de tension ne dépend que du nombre de spire.

On a $U_1 = R_1i_1 + \frac{\dd \phi_1}{\dd t}\simeq \frac{\dd \phi_1}{\dd t}$ et $U_2 = R_2i_2 + \frac{\dd \phi_2}{\dd t}\simeq \frac{\dd \phi_2}{\dd t}$ en négligeant la résistance.

On obtient $u_1 = L_1\frac{\dd i_1}{\dd t} + M\frac{\dd i_2}{\dd t}$ et
$u_2 = L_2\frac{\dd i_2}{\dd t} + M\frac{\dd i_1}{\dd t}$

On isole $\frac{\dd i_1}{\dd t} = \frac{1}{L_1}(u_1-M\frac{\dd i_2}{\dd t}) = \frac{1}{M}(u_2-L_2\frac{\dd i_2}{\dd t})$

On réécrit $u_2 = \frac{M}{L_2}u_1 + \frac{L_1L_2-M^2}{L_1}\frac{\dd i_2}{\dd t}$

On note $k$ le coefficient de couplage : $k = \frac{M}{\sqrt{L_1L_2}}$. Dans le cas idéal, k=1 pour avoir tout le flux d'une bobine qui
passe dans l'autre. Dans ce cas 
$$u_2 = \frac{M}{L_1}u_1$$ et un flux total $\phi$ qui est le m\^eme.

On a donc $L_1 = \frac{N_1\phi}{i_1}$ et $M=\frac{N_2\phi}{i_1}$

On en déduit que $\frac{M}{L_1}=\frac{N_2}{N_1}$ et finalement $\frac{u_2}{u_1} = \frac{N_2}{N_1}$


\end{document}
